\section{Core Mechanisms}
\label{sec:mech}
\begin{figure*}[!t]
\centering
\includegraphics[width=\textwidth]{figures/publication/fig3-channel-loop.pdf}
\caption{Channel organization and supervisory feedback. A lead coordinates member instances through tasks, messages, and memory; node bindings separately delimit industrial actions. Proposals pass through source-appropriate checks to device execution and observations. Agent judgment records a conclusion, whereas the independent backstop can initiate a separate recovery action.}
\label{fig:teamloop}
\end{figure*}
The six core mechanisms and the cross-cutting HITL boundary below connect commissioning, collaboration, execution, production context, feedback, recovery, and human intervention. Fig.~\ref{fig:teamloop} introduces their organizational setting; Section~\ref{sec:loop} details the feedback mechanism rather than equating the channel's task loop with physical process control.

\subsection{Node Connection and Creation}
\label{sec:nodes}
A node separates the meaning of an industrial variable from its transport address. Commissioning associates device integration information with acquisition or control configuration, tests connectivity, and supplies units, calibration, limits, and production-line membership. Drivers support Modbus TCP, Modbus RTU over TCP, OPC UA, MQTT, and HTTP/REST; a mock implementation supports testing. Device associations are many-to-many metadata: a device may be associated with several nodes, and a node may reference several devices. These associations do not imply that one control-node write fans out across multiple connections. Consequently, device identity, node identity, and agent identity are distinct rather than interchangeable keys.

Creation also differs from authorization to operate. Operator APIs and agent-facing creation paths have different permission conditions; the architecture does not establish identical authorization at every ingress. For industrial agent tools, bindings refer to concrete channel-member instances, not reusable agent templates. Deploying the same template elsewhere therefore does not, by itself, transfer a node binding. Team organization can simplify assignment, but membership alone is not write permission.

The semantic card presents the node's physical quantity, unit, precision, range, current value, freshness, and available recipe context. It gives reasoning a process-level vocabulary while keeping protocol encoding within the driver layer. The recommendation to change a setpoint by at most 2\,\% of its span is prompt guidance, not a controller-enforced step limit. Per-node acquisition runtimes separately manage due times and in-flight sampling. A rotating gateway sweep apportions sampling work, while acquisition, publication, and display intervals serve different purposes. This scheduling supports supervisory observation; it does not establish a hard real-time sampling guarantee.

\subsection{The Agent Channel: Team, Supervision, and Memory}
\label{sec:channel}
A channel contains a lead, worker instances, a task tree, and message history. Routing resolves recipients within the channel, loads runtimes when necessary, and delivers messages through mailboxes or supported steering interfaces. The task tree distinguishes submitted, assigned, working, waiting, failed, cancelled, and completed work. These states describe collaboration progress, not device execution state: completing an analysis task is neither permission to write nor evidence of a successful physical intervention.

The scheduler combines periodic supervision with event wake-up and idle backoff. It summarizes tasks, queues, member activity, and recent messages for the lead, avoiding redundant supervision when the relevant state is unchanged. Deterministic rules handle dispatch, retry, and stalled work; recent tool activity contributes to liveness detection. Model reasoning can therefore choose a course of action without being responsible for every scheduling transition. Conversely, scheduler progress cannot resolve uncertainty about a plant response; that requires observations and an explicit judgment.

Engine adapters support interchangeable reasoning implementations, including process-external engines. Channel memory preserves task experience, reflections, and reusable knowledge for later retrieval. Its function is informational continuity: a subsequent worker can recover the rationale and outcome of earlier work rather than infer them from a current setpoint alone. Retrieved memories remain contextual evidence, not an alternative authorization mechanism or an executable safety specification. Industrial tools still resolve the invoking member's bindings when an action is requested.

\subsection{Node Control Execution}
\label{sec:writepath}
The execution path separates an authorized request from a verified device result. An ordinary agent setpoint request first resolves a control-node binding and, in manual mode, waits for approval. The controller rejects requests when its gateway or the node is paused. Agent-origin requests additionally encounter optimization ownership and directional cooldown checks. The node runtime validates the engineering range and prevents overlapping writes before invoking protocol execution.

Let $S_n=[l_n,u_n]$ be the node range, $W_{r,n}$ the parameter window found through the active line run's recipe, and $q$ the optional recipe-run identifier supplied with a request. Let $P_r$ contain the nodes with a parameter in that recipe. The effective numerical constraint at controller entry is
\begin{equation}
C(n,r,q)=
\begin{cases}
S_n\cap W_{r,n}, & r\neq\varnothing,\ q=\varnothing,\ n\in P_r,\\
S_n, & \text{otherwise}.
\end{cases}
\label{eq:constraint}
\end{equation}
An unspecified window endpoint is unbounded. This formula describes normal controller requests, not every invocation of the low-level executor. In particular, recipe application supplies $q$ and bypasses the active-window branch. Hold heartbeats share write exclusion but call the executor directly, without re-entering all controller checks. Ownership and cooldown checks likewise depend on source metadata rather than applying universally.

Calibration maps engineering values to the device domain and maps readback into physical units. The runtime's engineering tolerance is
\begin{equation}
\tau_n=\max\!\left(0.5\,10^{-d_n},\;0.005(u_n-l_n)\right),
\qquad |v_{\rm rb}-v|\leq\tau_n,
\label{eq:tolerance}
\end{equation}
where $d_n$ is display precision. For linear calibration, the executor converts this tolerance using the absolute scale before protocol comparison. A successful readback supports agreement with the commanded value, not proof that the process has reached its desired operating condition. Communication failure or readback mismatch must therefore remain distinct from subsequent process-effect judgment. If a driver provides no readback, displaying the commanded value as a fallback does not establish measured agreement.

\begin{algorithm}[!t]
\caption{Ordinary bound-agent setpoint request}
\label{alg:write}
\begin{algorithmic}[1]
\Require Member $a$, node $n$, engineering target $v$
\State $b\gets\textsc{FindBinding}(a,n)$; reject if absent
\If{$b.\mathrm{mode}=\mathrm{manual}$}
  \State Await approval; reject on denial or timeout
  \State Recheck binding; reject if removed
\EndIf
\State Reject if gateway paused or node disabled
\State Check agent ownership and directional cooldown
\State Reject if $v\notin C(n,r,\varnothing)$
\State Validate range; acquire node write exclusion
\State Await connection, protocol write, and readback
\State Release exclusion; preserve execution outcome
\If{execution reports success}
  \State Attempt anchor, record, and operations bookkeeping
\EndIf
\State \Return execution result and available record identifiers
\end{algorithmic}
\end{algorithm}

Algorithm~\ref{alg:write} summarizes this specific ingress. Asynchronous connection and protocol queuing separate checks from device action, so neither co-residency nor node write exclusion supplies check--act atomicity against changing recipe state. Successful writes trigger journal anchors and, for relevant sources, optimization records. Bookkeeping errors are caught without reversing the write result. Attribution is therefore an important traceability mechanism, but not an atomic commitment spanning storage and the PLC.

\subsection{Production-Line Operations Management}
\label{sec:lineops}
Line operations associate products, recipe definitions, and runs with the nodes that generate measurements or accept setpoints. The run provides a batch identity for historical interpretation: a temperature observed during one product's recipe should not silently become evidence about another product's operating window. Runtime associations support tagged queries, operational histories, and the selection of evidence for optimization records.

Starting a line validates its recipe and product, creates a run, and installs that run as active before writing its parameters individually. An active run is consequently not a certificate that all parameters were successfully applied. Individual outcomes must be examined, especially when communication fails partway through initialization. Acquisition is gated by active batches, so opening a run also establishes the context in which its line's measurements are collected. Stopping the line clears this context, closes up to 500 open optimization records through the current stop hook, and requests acquisition-offline status. Complete record closure beyond that limit is not guaranteed. It is a supervisory batch operation, not a PLC emergency stop.

Monitoring combines live values, node conditions, alarms, and operational events. Historical measurements and frames complement current-state displays, allowing operators and agents to inspect a response over time rather than interpret a single value in isolation. The device-twin scene renders these states and associated assets. Its renderer does not issue protocol commands, although Town controls can submit writes through APIs. Thus neither the scene's visual appearance nor the stopped state of a batch should be mistaken for a physical safety barrier. Between batches, acquisition-based coverage is limited by the run gate; plant-side protection must remain independent.

\subsection{Agent Closed-Loop Control}
\label{sec:loop}
The feedback sequence is observation, hypothesis, intervention, judgment, and retention or recovery. The agent inspects bound-node context and acquisition history, states an expected effect, and proposes an engineering setpoint. After a successful value-changing write, the rollback manager creates an anchor and an optimization record containing the node, actor, hypothesis, parameter change, and applicable policy. Baseline statistics are filled asynchronously from preceding observations. After a five-second same-target deduplication check, a subsequent successful write supersedes the node's open record. An unchanged-value write then creates no replacement anchor or record; a value-changing agent or rollback write starts a new observation record. This is node-scoped bookkeeping, not one experiment per agent across the plant.

Before an agent write, an open record belonging to another agent blocks entry unless it has become stale and is superseded. A recent rollback also prevents an agent from immediately moving in the same direction away from the restored value. These source-specific checks reduce conflicting interventions; they do not exclude manual or recipe-origin writes. Defaults include a 600\,s baseline lookback, a 300\,s cooldown, and a 30\,min stale-record interval, with configuration controlling the corresponding durations.

The agent then invokes \texttt{dcw\_judge} with a record identifier, a verdict, and a reason. The tool validates and records the supplied conclusion; it does not calculate that conclusion from telemetry. A \emph{keep}, \emph{rollback}, or \emph{uncertain} judgment is therefore attributable reasoning, not an independently established process result. Judgment itself performs no actuation. A keep judgment closes the record and can mark a good state, whereas rollback and uncertain judgments leave it open for further handling. These transitions make the record a coordination object linking the intervention to its interpretation, rather than a substitute for measurement. Restoring a setpoint requires a separate rollback operation, and uncertainty must not be interpreted as successful optimization.

An independent backstop periodically evaluates eligible open records. With default settings it waits at least 120\,s and rechecks at 30\,s intervals. Acquisition channels matching the control node's primary-device alias are preferred (the selector does not traverse every device association); otherwise selection falls back to the line. Historical query results are compared with configured acquisition monitoring windows, and at least three accumulated breaches trigger a rollback judgment. Missing monitoring windows do not become control-node hard-range tests. The queried observations, aggregation, and data availability constrain this evidence; the intervals alone do not prove a worst-case physical recovery deadline.

For an automatic-rollback policy, the backstop executes recovery only when the node's count of records marked \texttt{rolled-back} is below $K=2$. This count spans retained records for that node, not merely the current run, and includes recorded rollbacks irrespective of initiator. Because repository retention is bounded, it is not a permanent node-lifetime counter. Nor is $K$ a universal limit on every explicit rollback route. At the threshold, or under an approval-required policy, the backstop records an escalation instead of automatically writing. The distinction confines the bound to the implemented backstop decision rather than claiming globally bounded agent behavior.

\begin{figure*}[!t]
\centering
\includegraphics[width=\textwidth]{figures/publication/fig4-workflow.pdf}
\caption{Operator-facing workflow using archived interface crops: node commissioning, device association, team deployment, and optimization-record inspection. Device association and agent capability binding are distinct operations. A displayed judgment records a disposition; it does not independently certify process improvement or physical recovery.}
\label{fig:walkthrough}
\end{figure*}

\subsection{Recipe Management and Rollback}
\label{sec:recipe}
Recipe management separates the definition of intended operation from the device state actually achieved. Versioned edits and known-good snapshots preserve candidate parameter sets and their provenance. Reverting a definition creates a new version with restored metadata; it does not, by itself, restore PLC values. Recipe application instead creates a run record and dispatches its parameters, while line start additionally establishes the active line context. These operations must not be conflated merely because each refers to the same recipe. During parameter dispatch, deleted or disabled nodes and nodes no longer attached to the recipe's line produce explicit unsuccessful outcomes; other parameters can continue. The run retains these individual results, allowing partial application to be distinguished from complete application. This is essential when deciding whether to retry, correct the configuration, or compensate for parameters already changed.

The ordinary active-window check resolves the recipe through the active run's recipe identifier; it is not an atomic check against every frozen run field. Application passes its run identifier into controller writes and therefore skips that particular window branch, while retaining runtime range validation. These details matter when interpreting whether a version change alters the constraints encountered by a later request. Version history documents intent, but does not establish transactional consistency between an edited definition, all run metadata, and physical parameters.

Physical rollback selects a recorded baseline or another supported anchor, batch-start, or known-good target and issues a new write. The controller can reject that recovery, for example when paused or when the target violates an applicable current constraint. A rollback verdict therefore does not guarantee restoration. Across multiple nodes, recovery is compensating work: one restoration may succeed while another fails. There is no distributed physical transaction that erases intermediate process states. Operators must distinguish restored recipe metadata, successful individual compensation, and verified recovery of the production process.

\subsection{Human-in-the-Loop Control}
\label{sec:hitl}
Human intervention is attached to concrete bindings and operations. For an ordinary bound-agent control request, \emph{manual} mode suspends execution for approval; \emph{auto} mode omits that human wait but retains applicable controller and runtime checks. An unbound member is rejected. On resumption after approval, the tool verifies that the binding still exists before entering the controller. Approval thus permits a request to proceed, rather than reserving a device state or guaranteeing that later checks will pass.

Tool approvals have an explicit expiry timer and resolve as denied when unanswered. Duplicate pending requests for the same agent and node are suppressed, and removing a binding can cancel outstanding requests. This timed approval object must be distinguished from an optimization record carrying a rollback judgment awaiting human action. The latter is record-based state; it is not evidence that tool approvals lack a timeout, nor does a recorded escalation mean that recovery has executed.

Human decisions, agent proposals, and system recovery events provide complementary attribution. Their usefulness depends on retaining the distinction between intent, authorization, device outcome, and observed process effect. Process-external engines also require deployment-level privilege and network restrictions: tool bindings govern sanctioned tool calls, not arbitrary I/O by a compromised local process. Fig.~\ref{fig:walkthrough} links the mechanisms to four operator-facing surfaces: node commissioning, device association, team deployment, and optimization records. These archived interface crops demonstrate the available interaction surfaces, not a new successful control experiment or a study of human approval behavior.
